\section{Actual affine pullback germs and literal localization tensors}

\begin{lemma}[The original affine section module has its exact original scalar action]
\label{lem:affine-identity-sections-module}
\leanok
\lean{EnvelopeProperty.AffinePullbackStalkTensorReadback.identitySectionsEquiv,EnvelopeProperty.AffinePullbackStalkTensorReadback.identitySectionsEquiv_apply,EnvelopeProperty.AffinePullbackStalkTensorReadback.identitySectionsTensorEquiv,EnvelopeProperty.AffinePullbackStalkTensorReadback.identitySectionsTensorEquiv_tmul,EnvelopeProperty.AffinePullbackStalkTensorReadback.originalGerm_eq_zero_iff_identitySectionsTensor}
The module of original affine sections supplied by the spectrum global-section
functor is canonically linearly equivalent to the original section module
with coefficient map of the identity spectrum morphism. This equivalence
retains every original section literally, and so does its tensor extension. Consequently the actual original affine
germ is zero if and only if its tensor in the actual point-local ring,
using the identity-map section module, is zero.
\end{lemma}
\begin{proof}
Both modules have the same original additive group of global sections.
The first action of $r$ is through the inverse canonical spectrum/global
section isomorphism. The second action is through that same inverse
followed by the identity morphism's top section map, which is the identity.
Their scalar actions agree. The identity additive equivalence is therefore
linear and its tensor extension is the identity on every pure tensor.
The already proved actual affine germ/localization-tensor criterion
therefore transports through this SAME equivalence, without unfolding
the original sheaf or supplying a stalk-comparison certificate.
\end{proof}

\begin{lemma}[Literal original coefficients in the actual affine pullback tensor]
\label{lem:affine-pullback-sections-tensor-readback}
\leanok
\lean{EnvelopeProperty.AffinePullbackStalkTensorReadback.specPullbackTensorEquiv,EnvelopeProperty.AffinePullbackStalkTensorReadback.specPullbackTensorEquiv_tmul,EnvelopeProperty.AffinePullbackStalkTensorReadback.extendedPullbackTensorEquiv,EnvelopeProperty.AffinePullbackStalkTensorReadback.extendedPullbackTensorEquiv_tmul}
\uses{lem:affine-identity-sections-module}
For an original ring map $B\to T$ and original quasi-coherent module
$M$ on $\operatorname{Spec}B$, put $N=g^*M$ for its actual spectrum
map $g$. There is a canonical $T$-linear equivalence
$T\otimes_B\Gamma(M)\simeq\Gamma(N)$, carrying EVERY $t\otimes m$
to $t\,u(m)$ with the actual original adjunction-unit section $u(m)$.
For every original $T$-algebra $D$, its scalar extension and canonical
tensor cancellation give
$D\otimes_T\Gamma(N)\simeq D\otimes_B\Gamma(M)$, carrying
$d\otimes(t\,u(m))$ to $(t d)\otimes m$ with the literal original
composite $B\to T\to D$ action.
\end{lemma}
\begin{proof}
Use the canonical spectrum/global-section ring isomorphisms on both
original rings. Their naturality identifies the original $B\to T$
action with the actual top-section map. The identity additive map on
$\Gamma(M)$ is semilinear for the original spectrum coefficient ring
isomorphism, and the target spectrum coefficient isomorphism is
semilinear for the same base isomorphism. Their tensor equivalence,
followed by the actual affine pullback global-section equivalence,
gives the asserted additive equivalence. Multiplication of every pure
coefficient tensor and additivity prove $T$-linearity on all tensors.
The actual affine pullback equivalence sends $1\otimes m$ to the
original adjunction unit. Linearity gives the stated formula for every
$t$. Extend this genuine equivalence to $D$ and cancel the intermediate
$T$ tensor by canonical associativity and multiplication; its pure
tensor formula is exactly $(td)\otimes m$.
\end{proof}

\begin{theorem}[Actual pulled-back section germs are the literal original localized tensors]
\label{thm:affine-pullback-stalk-tensor-readback}
\leanok
\lean{EnvelopeProperty.AffinePullbackStalkTensorReadback.pullbackGerm_eq_zero_iff_localizedTensor,EnvelopeProperty.AffinePullbackStalkTensorReadback.pullbackGerm_eq_zero_iff_identitySectionsTensor}
\uses{lem:affine-identity-sections-module,lem:affine-pullback-sections-tensor-readback}
Let $B\to T$ be an original ring map, $g:\operatorname{Spec}T\to
\operatorname{Spec}B$ its actual spectrum morphism, and $M$ an original
quasi-coherent module sheaf on $\operatorname{Spec}B$. Write $N=g^*M$,
and let $u(m)$ be the actual top section supplied by the original
pullback-adjunction unit. For a prime $x$ of $T$, put $D=T_x$.
For every original $m\in\Gamma(M)$ and EVERY $t\in T$,
\[
 (t\,u(m))_x=0\quad\Longleftrightarrow\quad
 (t/1)\otimes_B m=0\ \text{ in }D\otimes_B\Gamma(M).
\]
No stalk-comparison certificate and no flatness hypothesis is supplied.
The same statement consequently applies to the actual affine coefficient
tensor chart $T=B\otimes_R S$ and $t=1\otimes s$ for every $s\in S$.
\end{theorem}
\begin{proof}
Apply the preceding extended equivalence with the actual local ring
$D=T_x$. Its literal formula sends $1\otimes(t\,u(m))$ to
$(t/1)\otimes m$, so these tensors vanish together. The actual pullback
$N$ is quasi-coherent on the affine spectrum. The genuine affine
localization criterion therefore identifies vanishing of the original
$N$-germ with vanishing of its tensor $1\otimes(t\,u(m))$.
Combining these two equivalences proves the assertion. The first lemma
normalizes to the same identity-map section module when required.
\end{proof}
